\begin{prob}
    For each of the following pieces of code, state the time complexity using
    $\Theta$ notation in as simple of terms as possible. You do not need to
    show your work (but be careful, because without work we can't give you
    partial credit!).

    \begin{subprobset}

        \begin{subprob}
            \inputminted{python}{\thisdir/include/part_a.py}

            \begin{soln}
                $\Theta(n^2)$.

                This is very similar to the dependent loops we saw in the
                tallest doctor problem where we considered unordered pairs.
            \end{soln}
        \end{subprob}

        \begin{subprob}
            \inputminted{python}{\thisdir/include/part_b.py}

            \begin{soln}
                $\Theta(n^6)$.

                Although we have a \python{while}-loop, it essentially
                behaves like a nested \python{for}-loop over \python{range(n)}.
                So we have independent loops, and the total number of executions
                of the innermost loop body is $\Theta(n^6)$.
            \end{soln}
        \end{subprob}

        \begin{subprob}
            \inputminted{python}{\thisdir/include/part_c.py}

            \begin{soln}
                $\Theta(n^2)$.

                When analyzing the time complexity of a piece of code where
                each line takes constant time to execute once, we just need to find the line
                that executes the most number of times. In this case, it's the
                inner body of the nested \python{for}-loop. It executes $n^2$ times, for
                a time complexity of $\Theta(n^2)$.

                We don't know if the last $\python{for}$ loop will run or not
                -- that depends on whether $n$ is even or odd. But even if it's
                even (no pun intended) and it \emph{does} run, it won't change
                the overall time complexity because $n^2 + n$ is still
                $\Theta(n^2)$.
            \end{soln}
        \end{subprob}

        \begin{subprob}
            \inputminted{python}{\thisdir/include/part_d.py}

            \begin{soln}
                $\Theta(n^2)$.

                This is an example of a piece of code where not every line
                takes constant time to execute! In particular, the \python{sum{arr}} line
                takes $\Theta(n)$ time to execute once, and it is executed $n$ times, for
                a total time of $\Theta(n^2)$.
            \end{soln}
        \end{subprob}

    \end{subprobset}
    
\end{prob}
